% Dimensional Game Theory — NFA DAG HDIS LaTeX
% File: DGT_NFA_HDIS.tex
% Author: OBINexus Research Group (Nnamdi Michael Okpala)
% Repo: github.com/obinexus/dgt
% Date: October 2025

\documentclass[12pt,a4paper]{article}
\usepackage{textgreek}
\usepackage{graphicx}
\usepackage{tikz}
\usepackage{tikz-cd}
\usepackage{hyperref}
\usepackage{geometry}
\usepackage{booktabs}
\usepackage{longtable}
\usepackage{enumitem}
\usepackage{caption}
\usepackage{amssymb}
\usepackage{listings}
\usepackage{algorithm}
\usepackage{algpseudocode}
\usepackage{float}
\geometry{margin=1in}
\hypersetup{colorlinks=true,linkcolor=blue,citecolor=blue,urlcolor=blue}

\title{\LARGE\bfseries Dimensional Game Theory \\
Application of Non-Deterministic Finite Automaton Directed Acyclic Graph \\
for Actor Modelling via Phenomenological Observer and Consumer \\
HDIS Functor Coherence Operator Interpolation with Fault-Tolerant Resolution Implementation}
\author{OBINexus Research Group \\
\texttt{github.com/obinexus/dgt}}
\date{October 2025}

\begin{document}
\maketitle
\begin{abstract}
This document formalises the Dimensional Game Theory (DGT) framework and its
Non-Deterministic Finite Automaton (NFA) functor extension, integrates
Directed Acyclic Graph (DAG) topology for HDIS (Hybrid Directed Instruction
Systems), and provides algorithmic and verification tooling (ODTS, CRIF).
It is intended as a single-file LaTeX source suitable for Overleaf compilation
and inclusion in the project repository. The text contains formal definitions,
algorithms, pseudocode, diagrams, and a practical integrity metric (CRIF) for
witness ledger verification.
\end{abstract}

\tableofcontents
\newpage

% --------------------
\section{Introduction}
\subsection{Background: From Classical Game Theory to Dimensional Extensions}
A short motivation for DGT: traditional game theory analyses strategic agents
with fixed action sets and payoffs. DGT extends this by adding ``strategic
dimensions'' that are activated contextually and by building functorial
mappings between observer and actor spaces.

\subsection{Role of Phenomenology in Computational Modelling}
Phenomenological elements (phenotype, phenomemory, phenovalue) are included
as first-class types to model human-in-the-loop systems and to ground
coherence metrics in measurable observables.

\subsection{HDIS as Foundation: Human-in-the-Loop Coherence Systems}
HDIS is introduced as the architectural seed for active systems that observe,
adapt, and self-heal. Implementations include DIRAM, DIPAD and the SemVerX
hot-swap engine.

\subsection{Research Aim and Scope}
This document unifies formalism, algorithms and reference implementations
(links to repositories) for DGT, NFA functor extension, ODTS and HDIS
integration.

% --------------------
\section{Theoretical Framework}
\subsection{Dimensional Game Theory (DGT): Core Definitions}
\subsubsection*{Glossary and Notation}
\begin{description}
  \item[$\mathbb{N}$] Natural numbers $\{0,1,2,\dots\}$.
  \item[$\mathbb{Z}$] Integers $\{\dots,-2,-1,0,1,2,\dots\}$.
  \item[$\mathbb{Q}$], $\mathbb{R}$, $\mathbb{C}$ standard sets of rational,
  real and complex numbers.
\end{description}

\subsubsection*{DGT-specific symbols}
\begin{description}
  \item[$F,G$] Functorial mappings between behaviour spaces.
  \item[$F\cdot G$] Composition: $(F\cdot G)(x)=F(G(x))$; in DGT composition may
    be non-deterministic through functor-lifting to power-set codomains.
  \item[$F\star G$] Coherence operator (lossless interaction), defined in
    Section~\ref{sec:coherence-operators}.
  \item[$\Phi,\Psi$] Phenomenological functors: $\Phi$ maps raw observations to
    phenovalue; $\Psi$ maps intent/actor transforms.
\end{description}

\subsection{Non-Deterministic Finite Automata (NFA) as Behavioural Models}
Define an actor as $\mathcal{A}_{\mathrm{NFA}}=(Q,\Sigma,\delta,q_0,F)$ where
$\delta:Q\times\Sigma\to 2^Q$. Actor policies are modelled as NFA transitions
and functor lifts let us compose NFAs with DAG topologies.

\subsection{Directed Acyclic Graphs (DAG) in System Coherence}
Nodes model actor components; edges model observer-consumer transactions and
coherence receipts. A DAG $G=(V,E)$ is used so that topological traversal
and functor composition are well-defined.

\subsection{Actor-Observer-Consumer Paradigm}
Each node implements an \emph{Observer} that produces phenovalue and a
\emph{Consumer} that acts on it. The pair yields dual receipts stored in the
witness ledger.

\subsection{Phenomenological Type System (Phenotype, Phenomemory, Phenovalue)}
Formal type definitions and interfaces for phenomodels, including minimal
API examples (pseudocode) for observing and consuming layers.

\subsection{Coherence Operators and Functor Composition}
\label{sec:coherence-operators}
Define algebraic operators used to combine functors: union $\mathcal{U}$,
coherence $\star$, disjoint $\div$, and null $\bot$. A working definition:
\[ H = F\star G = \mathcal{C}(F,G)\cdot (F\cdot G) \]
with $\mathcal{C}(F,G)\in[0,1]$ the coherence measure (see Section~\ref{sec:crif}).

% --------------------
\section{Mathematical and Algorithmic Foundation}
\subsection{State Transition Functions and NFA Functor Definition}
Formalisation of lifted transitions: if $\delta$ is the NFA transition,
we define the functor-lifted transition $\delta^\sharp$ mapping distributions
or power-set states.

\subsection{Functor Composition: \texorpdfstring{$F \cdot G$}{F.G} Operator Interpolation}
Describe the interpolation semantics when $F$ and $G$ are stochastic/non-
deterministic functors. Provide pseudocode for composition and for estimating
$\mathcal{C}(F,G)$ via sampled traces.

\subsection{Set-Theoretic Mapping: Union, Disjoint, Pairing, Division}
Formal set operations used to reason about lossless vs lossy joins in the DAG.

\subsection{Complexity Constraints and \texorpdfstring{$O(\log n)$}{O(log n)} Auxiliary Space}
State the Functor Framework's constraint: implement traversals and QA
aggregations using $O(\log n)$ auxiliary space (practical implementation
notes for iterative topological traversal with bounded frontier are provided).

\subsection{Lossless vs Lossy Systems (Huffman--AVL Equilibrium)}
Explain the equilibrium idea: Huffman-like compression (entropy minimisation)
balanced with AVL-like structural constraints to preserve recoverability.

\subsection{Coherence Preservation under Non-Determinism}
Lossless and lossy examples and the probing functor $P(\mathcal{A}_{\mathrm{NFA}},\theta)$
that measures sensitivity $\nabla_\theta \mathcal{C}$.

% --------------------
\section{OBINexus Derivative Tracing System (ODTS)}
A safety-first subsystem to trace symbolic derivatives and confirm termination
for polynomial-type models used by controllers. Example use-case and API.

% --------------------
\section{HDIS Integration}
\subsection{HDIS as a Consumer-Observer Digital Nervous System}
Architectural overview, the witness ledger concept and coherence receipts.

\subsection{Mapping Observer and Consumer Roles}
Concrete mapping table and role responsibilities.

\subsection{SemVerX: Evolutionary Component Replacement}
Specification of SemVerX, hot-swap policies and DAG-based resolver.

\subsection{Witness Ledger and Coherence Receipts}
Full CRIF specification and pseudocode (see Section~\ref{sec:crif}).

\subsection{Polyglot Implementation: Rust -- ESP32 -- Python}
Integration notes, canonical IR, and polyglot binding rules for the HDIS stack.

% --------------------
\section{Coherence Receipt Integrity Function (CRIF)}
\label{sec:crif}
% Insert the CRIF material provided earlier (equations, algorithm and defaults)

\subsection{Notation}
$G=(V,E)$ DAG, witness ledger $L=\{r_e\}$, receipts $r_e=(o,c,\phi,\tau,\sigma)$.

\subsection{CRIF Definition}
\[ C = \lambda\, I_{\mathrm{hash}}(L) + (1-\lambda)\, Q_{\mathrm{agg}} \]
with detailed computation steps and pseudocode implemented as Algorithm~\ref{alg:crif}.

\begin{algorithm}[H]
\caption{ComputeCRIF}
\label{alg:crif}
\begin{algorithmic}[1]
\Procedure{ComputeCRIF}{$G,L,Pipelines,w,\alpha,\lambda,C_{req}$}
\State compute $I_{hash}$ by validating hashes for receipts in $L$
\State compute per-pipeline $Q(P)$ using F1 gate metrics
\State compute $Q_{agg}=\sum_i \alpha_i Q(P_i)$
\State $C = \lambda I_{hash} + (1-\lambda) Q_{agg}$
\State \textbf{return} $C$ and pass/fail flag $C \ge C_{req}$
\EndProcedure
\end{algorithmic}
\end{algorithm}

% --------------------
\section{Actor Behaviour Modelling}
\subsection{Non-Deterministic Action-State Space}
Formal NFA-based actor models and compositional rules.

\subsection{Dimensional Extension of Actor Roles}
How dimensions activate and interact; detection algorithm (PCA-based) for
finding dominant strategic dimensions.

\subsection{Real-World Example: Housing Entanglement Scenario}
A worked example mapping public services failures into DGT observables and
ledger evidence. This section shows how to structure FOI/appeal evidence as
phenovalue traces that the HDIS network can witness.

\subsection{Fault-Tolerant Resolution via Phenomenological Feedback}
Active repair flow and decision paths.

\subsection{Adaptive Decision Pathways}
Algorithmic implementation of adaptive response (Algorithm~\ref{alg:adaptive}).

\begin{algorithm}[H]
\caption{AdaptiveResponse}
\label{alg:adaptive}
\begin{algorithmic}[1]
\Procedure{AdaptiveResponse}{$g,\hat{s}_o,D$}
\State detect dominant $D_{dom}=\{D\mid E(\hat{s}_o,D)>\theta\}$
\For{each $D\in D_{dom}$}
  \State generate counter-strategy $s^\mathrm{c}_D$
\EndFor
\State combine counters weighted by dominance
\State \textbf{return} combined strategy
\EndProcedure
\end{algorithmic}
\end{algorithm}

% --------------------
\section{Simulation and Verification}
\subsection{DAG Construction and Traversal Pseudocode}
Topological traversal, bounded auxiliary stack and checkpointing notes.

\subsection{NFA-Functor Evaluation Loop}
Pseudocode for sampling NFA traces, lifting to functors and computing
$\mathcal{C}$ estimates.

\subsection{Huffman--AVL Compression for State Balancing}
Outline of combining Huffman coding heuristics with AVL balancing to reduce
state representation entropy while preserving recoverability.

\subsection{Complexity Validation Test Cases}
Example test harness and expected pass/fail criteria.

\subsection{Coherence Score Metrics (\textDelta-Meaning Scale)}
Definition of the witness score used in the registry (scale -12..+12 or 0..1
normalised). Recommended thresholds.

% --------------------
\section{Implementation Architecture}
\subsection{System Layers: IaaS, BaaS, Observer Interfaces}
Stack diagram and responsibilities.

\subsection{OpenSense Integration and Sensory Phenotyping}
Sensor mapping examples (IMU, gaze, skin-response) into phenovalue types.

\subsection{Registry and Witness NFT Economy}
How coherence receipts can be packaged as signed artefacts for funding
applications and audit trails.

\subsection{Fault-Tolerant Hot-Swap Engine (SemVerX Resolver)}
Resolver pseudocode with Eulerian/Hamiltonian path scoring.

\subsection{Prototype: \texttt{hdis-d} Daemon and ESP32 Module}
Quickstart: clone, build and flash steps; see README in repo.

% --------------------
\section{Discussion}
\subsection{Phenomenological Implications for Human Computation}
Ethical stakes, agency and the risk of over-automation.

\subsection{Socio-Technical Impact in Public Infrastructure}
How witness receipts and coherence scores could change accountability.

\subsection{Funding and Accessibility Mitigation via Witness Proofs}
Examples of grant/funding documentation derived from ledger evidence.

\subsection{Ethical and Legal Dimensions of HDIS}
Privacy, consent, and recommended governance model.

% --------------------
\section{Conclusion}
\subsection{Synthesis of DGT--NFA--Functor Model}
\subsection{Future Research Directions (DGT++ and Quantum Coherence)}
\subsection{Closing Reflection}

\section*{References}
% Minimal placeholder references; expand before submission.
\begin{thebibliography}{9}
\bibitem{neumann44} J. von Neumann and O. Morgenstern, \emph{Theory of Games and Economic Behavior}, 1944.
\bibitem{nash50} J. Nash, \emph{Equilibrium points in n-person games}, PNAS, 1950.
\bibitem{obinexus} N. M. Okpala, \emph{Dimensional Game Theory: DGT (draft)}, OBINexus, 2025.
\end{thebibliography}

\end{document}
